В повторение:
Определитель матрицы с углом нулей.
Определитель треугольной матрицы.
Линейная независимость векторов


В основной материал:
Линейная независимость подпространств.
Direct sum of subspaces.
Linear operators on direct sums of invariant subspaces


В задачи: 


Из Тимашева
\subsection{Тимашев, лекция 8}

\begin{definition}
Корневой вектор и его высота.
\end{definition}

(Корневой вектор высоты 0, 1.
Примеры корневых векторов)

\begin{definition}
Корневое подпространство, соответствующее собственному значению $\lambda$.
\end{definition}

План: 
1) доказать, что пространство раскладывается в прямую сумму корневых подпространств. Тогда матрица оператора в базисе, согласованном с базисом корневых подпространств, будет блочно-диагональной.
2) изучить матрицу на каждом блоке.

(Если есть вектор высоты $k$, то все включения строгие и в конечномерном пространстве включения строгие и размерности растут, но однажды стабилизируются. Значит, высота не превосходит размерности пространства.)

\begin{lemma}
Корневое подпространство инвариантно относительно оператора. (И относительно $A-\lambda I$.)

$A-\lambda I$ на своем корневом подпространстве нильпотентно.

($A-\mu I$ невыржодено на чужом корневом подпространстве ($\mu \neq \lambda$).)

Размерность корневого подпространства равна алгебраической кратности собственного значения.
\end{lemma}

\begin{theorem}
    Корневые подпространства, соответствующие различным собственным значениям, линейно независимы.
\end{theorem}

\begin{corollary}
Получили искомое разложение в прямую сумму.
\end{corollary}
